% HAND-WRITTEN fragment -- NOT generated by maestro2tex. Input directly from
% AnalogChapter.tex, before the block sections, because every block refers to it.
%
% The component values below are the single source of truth for the electrochemical
% models. Do NOT restate them in figure captions: a caption copy of r_ct drifted to
% 10k against the bench's 30k and shipped in the TRM before this section existed.
% Values read from: castalia/randles_cell and castalia/randles_electrode (OA), and
% confirmed against the instance line in
% simarchive/anatop_tb_controlamp/cm_2026-07-27/Interactive.156/1/CellDriveTB/netlist/.

\subsection{Electrochemical Models}\label{ss:analog-models}

Every electrode measurement in this chapter uses one of the two lumped models of
Figures~\ref{fig:randles-cell} and~\ref{fig:randles-electrode}, valued in
Table~\ref{tab:analog-models}. The electrode impedance sets the feedback factor
of both the potentiostat and the transimpedance loop, so those blocks' stability
figures depend on it directly.

The three-electrode cell is what the potentiostat drives: into the working
electrode, \(R_{\mathrm{ct}} \parallel C_{\mathrm{dl}}\) in series with
\(R_u\), about \SI{31}{\kilo\ohm} at DC and falling toward
\(R_u = \SI{1}{\kilo\ohm}\) above roughly \SI{5}{\hertz}. No bench reported here
uses the single-interface model.

Both models are linear and frequency-independent, with no Warburg term and no
potential dependence, so they give small-signal behaviour about a bias point
rather than a full electrochemical response. Both are also nominal: no element
carries a spread, so nothing in this chapter accounts for electrode-to-electrode
variation, which on a real wound interface is likely to exceed the process
spread reported beside it.

\begin{table}[htbp]
  \centering
  \caption[{Component values of the electrochemical models.}]{Component values of
  the electrochemical models. The three-electrode values are the subcircuit
  defaults; the transimpedance benches override them, and the values used are
  stated with each measurement.}
  \label{tab:analog-models}
  \begin{tabular}{llrl}
    \toprule
    Model & Element & Value & Represents \\
    \midrule
    three-electrode cell      & $R_{\mathrm{ce}}$  & \SI{1}{\kilo\ohm}    & solution resistance, CE to RE \\
                                & $R_u$              & \SI{1}{\kilo\ohm}    & uncompensated resistance, RE to WE surface \\
                                & $R_{\mathrm{ct}}$  & \SI{30}{\kilo\ohm}   & charge transfer at the WE interface \\
                                & $C_{\mathrm{dl}}$  & \SI{1}{\micro\farad} & double layer at the WE interface \\
    \addlinespace
    single interface & $R_s$              & \SI{10}{\kilo\ohm}   & series solution resistance \\
                                & $R_{\mathrm{ct}}$  & \SI{1}{\mega\ohm}    & charge transfer \\
                                & $C_{\mathrm{dl}}$  & \SI{10}{\nano\farad} & double layer \\
    \bottomrule
  \end{tabular}
\end{table}

The closed-loop measurements of Section~\ref{ss:ampab} override the defaults with
$R_u = \SI{100}{\ohm}$ or \SI{1}{\kilo\ohm}, $R_{\mathrm{ct}} =
\SI{10}{\kilo\ohm}$ or \SI{1}{\mega\ohm}, and $C_{\mathrm{dl}} =
\SI{100}{\nano\farad}$ or \SI{1}{\micro\farad}; the value used is stated with
each measurement.
